\section{Примеры простейших систем массового обслуживания}

Математическим аппаратом для моделирования систем массового обслуживания (СМО) традиционно служат однородные марковские цепи с непрерывным временем. Эволюция такой цепи полностью определяется заданием начального распределения и матрицы переходных вероятностей $P(t)$. 

В случае сильной эргодичности цепи существует единственное стационарное распределение $\pi$, которое совпадает с предельным распределением вероятностей состояний и не зависит от начального состояния системы. Прямое вычисление элементов матрицы $P(t)$ через дифференциальные уравнения Колмогорова зачастую сопряжено со значительными аналитическими трудностями. На практике задача нахождения стационарных характеристик сводится к построению матрицы интенсивностей переходов $Q$ (инфинитезимального генератора) и последующему решению алгебраической системы $Q^T \pi = 0$.

\begin{example}{Двухканальная СМО с отказами}
Рассмотрим систему массового обслуживания, состоящую из двух параллельных каналов. На вход системы поступает пуассоновский поток заявок с интенсивностью $\lambda$. Времена обслуживания заявок в каналах независимы и распределены по показательному закону с параметрами $\mu_1$ и $\mu_2$ соответственно. Алгоритм маршрутизации задается следующими правилами:
\begin{itemize}
    \item Если оба канала свободны, заявка занимает один из них с равной вероятностью.
    \item Если свободен только один канал, поступающая заявка направляется в него.
    \item Если оба канала заняты, заявка получает отказ и покидает систему.
\end{itemize}
Требуется построить стохастический граф переходов и найти стационарное распределение вероятностей состояний системы.
\end{example}

\begin{eproof}
Определим фазовое пространство системы $S = \{0, 1, 2, 3\}$. Введем случайный процесс $X(t) \in S$, характеризующий состояние системы в момент времени $t$:
\begin{itemize}
    \item $0$ — в системе нет заявок (оба канала свободны);
    \item $1$ — занят только первый канал;
    \item $2$ — занят только второй канал;
    \item $3$ — заняты оба канала.
\end{itemize}
Предполагается, что траектории процесса $X(t)$ являются непрерывными справа функциями времени.

\begin{center}
    \begin{tikzpicture}[
        scale=1.5,
        >={Stealth[length=3mm]},
        state/.style={circle, draw, minimum size=0.8cm, thick},
        every node/.style={font=\small}
    ]
        % Состояния
        \node[state] (s0) at (0,0) {0};
        \node[state] (s1) at (2,1.2) {1};
        \node[state] (s2) at (2,-1.2) {2};
        \node[state] (s3) at (4,0) {3};

        % Переходы из 0
        \draw[->, thick] (s0) to [bend left=15] node[midway, above left] {$\frac{\lambda}{2}$} (s1);
        \draw[->, thick] (s0) to[bend right=15] node[midway, below left] {$\frac{\lambda}{2}$} (s2);

        % Переходы обратно в 0
        \draw[->, thick] (s1) to [bend left=15] node[midway, below right] {$\mu_1$} (s0);
        \draw[->, thick] (s2) to [bend right=15] node[midway, above right] {$\mu_2$} (s0);

        % Переходы в 3
        \draw[->, thick] (s1) to[bend left=15] node[midway, above right] {$\lambda$} (s3);
        \draw[->, thick] (s2) to [bend right=15] node[midway, below right] {$\lambda$} (s3);

        % Переходы обратно из 3
        \draw[->, thick] (s3) to [bend left=15] node[midway, below left] {$\mu_2$} (s1);
        \draw[->, thick] (s3) to[bend right=15] node[midway, above left] {$\mu_1$} (s2);
    \end{tikzpicture}
    
    \vspace{2mm}
    \small\textit{Схема интенсивностей переходов двухканальной СМО.}
\end{center}

Вычислим элементы матрицы интенсивностей $Q = \|q_{ij}\|$. Рассмотрим поведение переходных вероятностей $p_{ij}(h)$ при $h \to 0+0$:

\begin{itemize}
    \item \textbf{Переходы из состояния 0:} За малое время $h$ в систему поступает одна заявка с вероятностью $\lambda h + o(h)$. Согласно условию, выбор свободного канала равновероятен. Следовательно, вероятности переходов равны $p_{01}(h) = \frac{1}{2}(\lambda h + o(h))$ и $p_{02}(h) = \frac{1}{2}(\lambda h + o(h))$. Переходя к пределу, находим $q_{01} = q_{02} = \frac{\lambda}{2}$. Вероятность одновременного поступления двух заявок есть $o(h)$, поэтому $q_{03} = 0$.
    \item \textbf{Переходы из состояния 1:} Система переходит в состояние 3, если поступает новая заявка (вероятность $\lambda h + o(h)$), а обслуживание текущей заявки не завершается (вероятность $1 - \mu_1 h + o(h)$). В силу независимости событий их совместная вероятность равна $(\lambda h + o(h))(1 - \mu_1 h + o(h)) = \lambda h + o(h)$, откуда $q_{13} = \lambda$. Переход в состояние 0 происходит при завершении обслуживания (вероятность $\mu_1 h + o(h)$) и отсутствии новых заявок ($1 - \lambda h + o(h)$), что дает $q_{10} = \mu_1$.
    \item \textbf{Переходы из состояния 2:} В силу симметрии рассуждений получаем $q_{23} = \lambda$ и $q_{20} = \mu_2$.
    \item \textbf{Переходы из состояния 3:} Система переполнена, новые заявки отбрасываются. Изменение состояния возможно лишь при завершении обслуживания. Если обработка заканчивается в первом канале, он освобождается, и система переходит в состояние 2: $q_{32} = \mu_1$. Аналогично при освобождении второго канала $q_{31} = \mu_2$.
\end{itemize}

Диагональные элементы определяются из свойства консервативности цепи $q_{ii} = -\sum_{j \neq i} q_{ij}$. Матрица интенсивностей принимает вид:
$$
Q = \mat{
    -\lambda & \frac{\lambda}{2} & \frac{\lambda}{2} & 0 \\
    \mu_1 & -(\mu_1 + \lambda) & 0 & \lambda \\
    \mu_2 & 0 & -(\mu_2 + \lambda) & \lambda \\
    0 & \mu_2 & \mu_1 & -(\mu_1 + \mu_2)
}
$$

Рассматриваемая цепь Маркова конечна и неразложима. По теореме об эргодичности конечных цепей она является сильно эргодической, что гарантирует существование единственного стационарного распределения $\pi = (\pi_0, \pi_1, \pi_2, \pi_3)^T$. Решим матричное уравнение $Q^T \pi = 0$, которое расписывается в систему балансных уравнений:
$$
\begin{cases}
    -\lambda \pi_0 + \mu_1 \pi_1 + \mu_2 \pi_2 = 0, \\
    \frac{\lambda}{2} \pi_0 - (\mu_1 + \lambda) \pi_1 + \mu_2 \pi_3 = 0, \\
    \frac{\lambda}{2} \pi_0 - (\mu_2 + \lambda) \pi_2 + \mu_1 \pi_3 = 0, \\
    \lambda \pi_1 + \lambda \pi_2 - (\mu_1 + \mu_2) \pi_3 = 0.
\end{cases}
$$

Выразим $\pi_1$ и $\pi_2$ через $\pi_3$. Вычитая третье уравнение из второго, получаем:
$$
(\lambda + \mu_1)\pi_1 - (\lambda + \mu_2)\pi_2 = (\mu_2 - \mu_1)\pi_3.
$$
Из четвертого уравнения следует $\lambda(\pi_1 + \pi_2) = (\mu_1 + \mu_2)\pi_3$. Выражая $\pi_2$ и подставляя в полученную разность, находим точные соотношения:
$$
\pi_1 = \frac{\mu_2}{\lambda}\pi_3, \quad \pi_2 = \frac{\mu_1}{\lambda}\pi_3.
$$
Подставляя найденное $\pi_1$ во второе уравнение системы, определяем $\pi_0$:
$$
\frac{\lambda}{2} \pi_0 - (\mu_1 + \lambda)\frac{\mu_2}{\lambda}\pi_3 + \mu_2 \pi_3 = 0 \implies \frac{\lambda}{2} \pi_0 = \frac{\mu_1\mu_2}{\lambda}\pi_3 \implies \pi_0 = \frac{2\mu_1\mu_2}{\lambda^2}\pi_3.
$$

Окончательное решение восстанавливается из условия нормировки $\pi_0 + \pi_1 + \pi_2 + \pi_3 = 1$. Вынеся $\pi_3$ за скобки, получаем:
$$
\pi_3 \left( \frac{2\mu_1\mu_2}{\lambda^2} + \frac{\mu_2}{\lambda} + \frac{\mu_1}{\lambda} + 1 \right) = 1.
$$
Обозначим константу нормировки как $C = \frac{2\mu_1\mu_2}{\lambda^2} + \frac{\mu_2}{\lambda} + \frac{\mu_1}{\lambda} + 1$. Тогда стационарные вероятности задаются вектором:
$$
\pi_0 = \frac{2\mu_1\mu_2}{C\lambda^2}, \quad 
\pi_1 = \frac{\mu_2}{C\lambda}, \quad 
\pi_2 = \frac{\mu_1}{C\lambda}, \quad 
\pi_3 = \frac{1}{C}.
$$
\end{eproof}

\begin{example}{Многоканальная СМО с ограниченной емкостью ($M|M|c|K$)}
Рассмотрим систему массового обслуживания, функционирующую по следующим правилам:
\begin{enumerate}
    \item Входящий поток заявок является пуассоновским с интенсивностью $\lambda$.
    \item Система содержит $c$ идентичных и независимых каналов обслуживания. Время обслуживания заявки в каждом канале распределено экспоненциально с параметром $\mu$.
    \item Общая емкость системы ограничена числом $K$ ($K \ge c$). Это означает, что в системе одновременно может находиться не более $K$ заявок (из них $c$ на обслуживании и $K-c$ в очереди).
    \item Дисциплина обслуживания: если все каналы заняты, новая заявка становится в очередь, если в ней есть свободные места. Если в системе уже находится $K$ заявок, вновь прибывшая заявка получает отказ и теряется.
\end{enumerate}

\textbf{Требуется:}
\begin{enumerate}
    \item Показав, что число заявок в системе $X(t)$ является процессом рождения и гибели, построить граф интенсивностей переходов.
    \item Найти стационарное распределение вероятностей состояний.
    \item Записать выражения для среднего числа заявок в системе и средней длины очереди.
\end{enumerate}
\end{example}

\begin{eproof}
\textbf{1. Анализ процесса и построение графа.}
Определим состояние системы $X(t)$ как число заявок, находящихся в системе в момент времени $t$. Фазовое пространство системы конечно: $S = \{0, 1, \dots, K\}$.
Рассмотрим переходы за малый промежуток времени $h$:

\begin{itemize}
    \item \textbf{Интенсивность рождения (поступления заявки):}
    Пусть в системе находится $i < K$ заявок. Вероятность поступления новой заявки за время $h$ равна $\lambda h + o(h)$. Поскольку поток ординарен, вероятность поступления двух и более заявок есть $o(h)$.
    $$
    p_{i, i+1}(h) = \lambda h + o(h) \implies \lambda_{i, i+1} = \lambda, \quad 0 \le i < K.
    $$
    Если $i=K$, новые заявки отбрасываются, переход в состояние $K+1$ невозможен ($\lambda_{K, K+1} = 0$).

    \item \textbf{Интенсивность гибели (обслуживания заявки):}
    Пусть в системе находится $i > 0$ заявок. Рассмотрим два случая:
    
    а) \textit{Неполная загрузка каналов} ($0 < i \le c$). Заняты $i$ каналов. Обслуживание закончится, если хотя бы один из $i$ каналов завершит работу. Вероятность завершения обслуживания в конкретном канале равна $\mu h + o(h)$. Вероятность того, что завершится обслуживание ровно в одном из $i$ каналов (сумма вероятностей несовместных событий):
    $$
    p_{i, i-1}(h) = C_i^1 (\mu h + o(h))(1 - \mu h + o(h))^{i-1} = i\mu h + o(h).
    $$
    б) \textit{Полная загрузка каналов} ($c < i \le K$). Работают все $c$ каналов (остальные $i-c$ заявок ждут в очереди). Интенсивность выхода определяется только работающими каналами:
    $$
    p_{i, i-1}(h) = C_c^1 (\mu h + o(h))(1 - \mu h + o(h))^{c-1} = c\mu h + o(h).
    $$
    
    Объединяя результаты, получаем выражение для интенсивности гибели $\mu_i$:
    $$
    \mu_i = \min\{i, c\} \cdot \mu = 
    \begin{cases} 
    i\mu, & 0 < i \le c, \\
    c\mu, & c < i \le K.
    \end{cases}
    $$
\end{itemize}

Поскольку переходы возможны только в соседние состояния, процесс $X(t)$ является процессом рождения и гибели. Построим граф переходов:

\begin{center}
\begin{tikzpicture}[
    scale=1.1,
    >=Stealth, 
    node distance=2.0cm, 
    state/.style={circle, draw, minimum size=1.0cm, thick},
    every node/.style={font=\small}
]
    % Состояния
    \node[state] (s0) {0};
    \node[state] (s1) [right of=s0] {1};
    \node[state] (s2) [right of=s1] {2};
    \node (dots1) [right of=s2] {$\dots$};
    \node[state] (sc) [right of=dots1] {$c$};
    \node (dots2) [right of=sc] {$\dots$};
    \node[state] (sK) [right of=dots2] {$K$};

    % Стрелки рождения (верх)
    \draw[->, thick] (s0) to[bend left=40] node[midway, above] {$\lambda$} (s1);
    \draw[->, thick] (s1) to[bend left=40] node[midway, above] {$\lambda$} (s2);
    \draw[->, thick] (s2) to[bend left=40] node[midway, above] {$\lambda$} (dots1);
    \draw[->, thick] (dots1) to[bend left=40] node[midway, above] {$\lambda$} (sc);
    \draw[->, thick] (sc) to[bend left=40] node[midway, above] {$\lambda$} (dots2);
    \draw[->, thick] (dots2) to[bend left=40] node[midway, above] {$\lambda$} (sK);

    % Стрелки гибели (низ)
    \draw[->, thick] (s1) to[bend left=40] node[midway, below] {$\mu$} (s0);
    \draw[->, thick] (s2) to[bend left=40] node[midway, below] {$\mathbf{2}\mu$} (s1);
    \draw[->, thick] (dots1) to[bend left=40] node[midway, below] {$\mathbf{3}\mu$} (s2);
    \draw[->, thick] (sc) to[bend left=40] node[midway, below] {$\mathbf{c}\mu$} (dots1);
    \draw[->, thick] (dots2) to[bend left=40] node[midway, below] {$\mathbf{c}\mu$} (sc);
    \draw[->, thick] (sK) to[bend left=40] node[midway, below] {$\mathbf{c}\mu$} (dots2);
\end{tikzpicture}
\end{center}

\textbf{2. Стационарное распределение.}
Цепь Маркова конечна и неразложима, следовательно, она является сильно эргодической и существует единственное стационарное распределение $\pi = (\pi_0, \dots, \pi_K)$.
Воспользуемся общей формулой для процессов рождения и гибели:
$$
\pi_j = \frac{\lambda_{j-1} \cdot \lambda_{j-2} \cdot \dots \cdot \lambda_0}{\mu_j \cdot \mu_{j-1} \cdot \dots \cdot \mu_1} \pi_0.
$$
Введем обозначение $r = \frac{\lambda}{\mu}$ (приведенная интенсивность потока). 

Для $j \le c$ в знаменателе стоит факториал $j!$:
$$
\pi_j = \frac{\lambda^j}{1\mu \cdot 2\mu \cdot \dots \cdot j\mu} \pi_0 = \frac{(\lambda/\mu)^j}{j!} \pi_0 = \frac{r^j}{j!} \pi_0.
$$
Для $j > c$ в знаменателе первые $c$ множителей дают $c!$, а остальные $j-c$ множителей равны $c$:
$$
\pi_j = \frac{\lambda^j}{c!\mu^c \cdot (c\mu)^{j-c}} \pi_0 = \frac{\lambda^j}{c! c^{j-c} \mu^j} \pi_0 = \frac{r^j}{c! c^{j-c}} \pi_0.
$$
Вероятность простоя $\pi_0$ находится из условия нормировки $\sum_{j=0}^K \pi_j = 1$:
$$
\pi_0 = \left( \sum_{j=0}^{c} \frac{r^j}{j!} + \frac{r^c}{c!} \sum_{j=c+1}^{K} \left(\frac{r}{c}\right)^{j-c} \right)^{-1}.
$$

\textbf{3. Характеристики эффективности.}
Математическое ожидание числа заявок в системе:
$$
\bar{n} = \sum_{n=0}^{K} n \cdot \pi_n.
$$
Средняя длина очереди (число заявок, ожидающих обслуживания):
$$
L_q = \sum_{n=c+1}^{K} (n-c) \cdot \pi_n.
$$
\end{eproof}

\begin{note}
В нотации Кендалла описанная выше модель обозначается как \textbf{$M|M|c|K$}. Полученные формулы являются классическими для теории телетрафика (формулы Эрланга для систем с ожиданием и потерями).
\end{note}

\section{Классификация систем массового обслуживания и нотация Кендалла}

Для построения корректной математической модели СМО необходимо формализовать основные параметры, определяющие её функционирование. Выделяют следующие ключевые характеристики:

\begin{enumerate}
    \item \textbf{Характеристики входящего потока:}
    \begin{itemize}
        \item \textit{Тип распределения интервалов:} детерминированный, экспоненциальный (марковский), общий и др.
        \item \textit{Свойства потока:} стационарность, ординарность, наличие последействия.
        \item \textit{Поведение заявок:} возможность ожидания в очереди, наличие ограничений на время ожидания (системы с нетерпеливыми заявками), поведение при переполнении (отказ или ожидание).
    \end{itemize}
    
    \item \textbf{Механизм обслуживания:}
    \begin{itemize}
        \item \textit{Закон распределения времени обслуживания.}
        \item \textit{Режим работы:} индивидуальная или групповая обработка.
        \item \textit{Зависимость от состояния:} меняется ли скорость обслуживания в зависимости от длины очереди.
    \end{itemize}
    
    \item \textbf{Структурная организация:}
    \begin{itemize}
        \item \textit{Число каналов обслуживания ($c$):} одноканальные и многоканальные системы.
        \item \textit{Конфигурация очередей:} общая очередь для всех каналов или отдельные очереди к каждому серверу.
        \item \textit{Емкость системы ($K$):} ограничение на суммарное число заявок (на обслуживании + в очереди).
    \end{itemize}
    
    \item \textbf{Дисциплина обслуживания (Queue Discipline):}
    \begin{itemize}
        \item \textbf{FCFS} (First Come, First Served) — в порядке поступления (FIFO).
        \item \textbf{LCFS} (Last Come, First Served) — последним пришел, первым обслужен (стек).
        \item \textbf{RSS} (Random Service Selection) — случайный выбор заявки из очереди.
        \item \textbf{PS} (Processor Sharing) — одновременное обслуживание всех заявок с разделением ресурса.
        \item \textbf{Priority} — наличие приоритетных классов (с вытеснением или без).
    \end{itemize}
\end{enumerate}

\begin{definition}{Нотация Кендалла}
Стандартная запись для кодирования параметров СМО имеет вид строки с разделителями:
$$ A \mid B \mid c \mid K \mid N \mid Z $$
где:
\begin{itemize}
    \item $A$ — распределение интервалов между поступлениями заявок;
    \item $B$ — распределение времени обслуживания;
    \item $c$ — число обслуживающих приборов (серверов);
    \item $K$ — емкость системы (по умолчанию $\infty$);
    \item $N$ — размер популяции источника заявок (по умолчанию $\infty$);
    \item $Z$ — дисциплина обслуживания (по умолчанию FCFS).
\end{itemize}
\end{definition}

\begin{note}
Для обозначения вероятностных распределений ($A$ и $B$) используются следующие сокращения:
\begin{itemize}
    \item $M$ (Markovian) — экспоненциальное распределение (пуассоновский поток);
    \item $D$ (Deterministic) — детерминированное (постоянное) время;
    \item $E_k$ — распределение Эрланга $k$-го порядка;
    \item $G$ (General) — произвольное распределение.
\end{itemize}
\end{note}

\begin{example}{Примеры нотации}
\begin{enumerate}
    \item $M|M|1$ — простейшая система: пуассоновский вход, экспоненциальное обслуживание, один канал, бесконечная очередь.
    \item $M|D|10$ — пуассоновский вход, фиксированное время обслуживания, 10 каналов.
    \item $M|G|1|\infty|PS$ — одноканальная система с разделением процессора и произвольным временем обслуживания.
\end{enumerate}
\end{example}

\section{Формула Литтла}

Рассмотрим произвольную систему массового обслуживания как «черный ящик», в который поступает поток заявок, и из которого выходят обработанные заявки.

\begin{center}
\begin{tikzpicture}[
    >=Stealth, 
    % Увеличиваем расстояние между узлами, чтобы тексту на стрелках хватило места
    node distance=4.5cm, 
    auto,
    block/.style={
        rectangle, 
        draw, 
        thick, 
        fill=Snow2, 
        text width=5.5cm, 
        align=center, 
        minimum height=1.5cm
    }
]
    % Центральный блок
    \node[block] (sys) {\textbf{Система} \\ (Очередь + Серверы)};
    
    % Точки входа и выхода (разнесены дальше)
    \node[left=of sys] (in) {};
    \node[right=of sys] (out) {};

    % Входящий поток: используем yshift, чтобы текст не касался линии и рамок
    \draw[->, ultra thick] (in) -- (sys.west) 
        node[midway, above=0.1cm] {$\lambda$} 
        node[midway, below=0.2cm, font=\small, text=gray] {Входящий поток};
        
    % Исходящий поток
    \draw[->, ultra thick] (sys.east) -- (out) 
        node[midway, below=0.2cm, font=\small, text=gray] {Обработанный поток};
    
    % Верхняя и нижняя метки
    \node[above=0.2cm of sys, text=red!70!black, font=\large] {$L$ (среднее число заявок)};
    \node[below=0.2cm of sys, text=cyan!50!black, font=\large] {$W$ (среднее время пребывания)};

    % Формула закона Литтла снизу
    \node[below=1cm of sys, font=\Large] {$L = \lambda W$};

\end{tikzpicture}
\end{center}

Введем формальные определения основных характеристик функционирования системы.
Пусть $A^{(k)}$ — момент поступления $k$-й заявки в систему ($k \in \mathbb{N}$), причем $0 \le A^{(1)} \le A^{(2)} \le \dots$.
Обозначим через $A(t)$ число заявок, поступивших в систему за интервал $[0, t)$.

Пусть $W^{(k)}$ — время пребывания $k$-й заявки в системе (время ожидания + время обслуживания). Тогда момент выхода $k$-й заявки: $D^{(k)} = A^{(k)} + W^{(k)}$.

Случайный процесс $N(t)$, характеризующий число заявок в системе в момент времени $t$, определяется как количество таких индексов $k$, для которых заявка уже поступила, но еще не покинула систему:
$$ N(t) = \sum_{k=1}^{\infty} \mathbb{I}\left( A^{(k)} \le t < A^{(k)} + W^{(k)} \right). $$

Предположим, что рассматриваемая система функционирует в установившемся (стационарном) режиме. Определим средние характеристики через пределы по времени и по числу заявок.

\begin{enumerate}
    \item \textbf{Интенсивность входного потока ($\lambda$):} предел отношения числа поступивших заявок к длительности интервала наблюдения:
    $$ \lambda \doteq \lim_{t \to \infty} \frac{A(t)}{t}. $$
    
    \item \textbf{Среднее время пребывания ($W$):} среднее арифметическое времен пребывания первых $k$ заявок при $k \to \infty$:
    $$ W \doteq \lim_{k \to \infty} \frac{1}{k} \sum_{i=1}^{k} W^{(i)}. $$
    
    \item \textbf{Среднее число заявок ($L$):} усредненное по времени значение процесса $N(t)$:
    $$ L \doteq \lim_{T \to \infty} \frac{1}{T} \int_{0}^{T} N(t) \, dt. $$
\end{enumerate}

\begin{theorem}{Закон Литтла}
Если для системы массового обслуживания существуют конечные пределы интенсивности входящего потока $\lambda$ и среднего времени пребывания заявки $W$, то существует и предел среднего числа заявок в системе $L$, причем справедливо равенство:
$$ L = \lambda \cdot W. $$
\end{theorem}

\begin{note}
Фундаментальная ценность закона Литтла заключается в его абсолютной универсальности. В формулировке теоремы не фигурируют ни законы распределения входящего потока и времени обслуживания, ни число каналов, ни дисциплина очереди (FCFS, LCFS, PS и др.). Тождество связывает исключительно асимптотические средние характеристики.
\end{note}

\begin{tproof}
Строгое аналитическое доказательство для систем с произвольной структурой сопряжено с громоздкими техническими деталями. Однако физический смысл закона прозрачно иллюстрируется геометрическим подходом. 

\textbf{1. Базовые допущения и геометрическая интерпретация.} 
Для начала предположим, что система функционирует по дисциплине FCFS (заявки покидают систему в порядке поступления), а в моменты времени $t=0$ и $t=T$ система гарантированно пуста. 

Вспомним, как мы ввели считающие процессы $A(t)$ и $D(t)$:
\begin{itemize}
    \item $A(t)$ — суммарное число заявок, поступивших в систему на интервале $[0, t)$;
    \item $D(t)$ — суммарное число заявок, покинувших систему на интервале $[0, t)$.
\end{itemize}
Тогда число заявок, находящихся в системе в произвольный момент времени $t$, определяется разностью $N(t) = A(t) - D(t)$. Пусть за время $T$ через систему прошло ровно $N$ заявок ($A(T) = D(T) = N$).

\begin{center}
\begin{tikzpicture}[>=Stealth, scale=1.3, font=\small]
    % Сетка и оси
    \draw[->, thick] (-0.5, 0) -- (9.5, 0) node[right] {$t$};
    \draw[->, thick] (0, -0.5) -- (0, 5.5) node[above] {Счетчик заявок};
    
    % Засечки на оси Y и горизонтальные пунктирные линии к ступеням A(t)
    \foreach \y/\x in {1/1, 2/3.5, 3/4, 4/4.5} {
        \draw[thin] (0.1, \y) -- (-0.1, \y) node[left] {\y};
        \draw[dashed, gray, thick] (0, \y) -- (\x, \y);
    }
    
    % Заливка области (Очередь)
    \fill[Snow2] 
        (1,0) -- (1,1) -- (3.5,1) -- (3.5,2) -- (4,2) -- (4,3) -- (4.5,3) -- (4.5,4) -- (8.5,4) -- 
        (8,4) -- (8,3) -- (7,3) -- (7,2) -- (6,2) -- (6,1) -- (2.5,1) -- (2.5,0) -- cycle;

    % График A(t) - Входящий поток (Красный)
    \draw[ultra thick, Red3] 
        (0,0) -- (1,0) -- (1,1) -- (3.5,1) -- (3.5,2) -- (4,2) -- (4,3) -- (4.5,3) -- (4.5,4) -- (8.5,4);
    \node[Red3, above left] at (4.5,4) {$A(t)$};

    % График D(t) - Выходящий поток (Синий)
    \draw[ultra thick, DeepSkyBlue4] 
        (0,0) -- (2.5,0) -- (2.5,1) -- (6,1) -- (6,2) -- (7,2) -- (7,3) -- (8,3) -- (8,4) -- (8.5,4);
    \node[DeepSkyBlue4, below right] at (6,1) {$D(t)$};

    % --- АННОТАЦИИ ВРЕМЕНИ ПРЕБЫВАНИЯ W ---
    \draw[<->, dashed] (1, 0.5) -- (2.5, 0.5) node[midway, above] {$W^{(1)}$};
    \draw[<->, dashed] (3.5, 1.5) -- (6, 1.5) node[midway, above] {$W^{(2)}$};
    \draw[<->, dashed] (4, 2.5) -- (7, 2.5); 
    % W(3) перенесен чуть ниже и правее точки пересечения
    \node[below right, inner sep=1pt] at (5.25, 2.5) {$W^{(3)}$};
    
    \draw[<->, dashed] (4.5, 3.5) -- (8, 3.5) node[midway, above] {$W^{(4)}$};

    % --- СПЕЦИАЛЬНЫЕ МЕТКИ ---

    % 1. Система пуста (середина графика)
    \draw[->, thick, black] (3, -0.65) -- (3, 0.95);
    \node[below, align=center, font=\footnotesize] at (3, -0.7) {Система\\простаивает\\($A=D$)};

    % 2. В системе 3 заявки (Зеленым цветом, как на картинке)
    % Вертикальная размерная линия
    \draw[<->, thick, EXAMP_COLOR] (5.2, 1.05) -- (5.2, 3.95);

    % Надпись N(t)=3 справа от линии, темно-зеленая
    \node[EXAMP_COLOR, anchor=west, font=\bfseries] at (5.2, 3.0) {$N(t)=3$};

    % Концевые метки
    \draw[dashed, thick] (8.5, 5.0) -- (8.5, 0) node[below] {$T$};
    \node[below, font=\footnotesize] at (0, -0.4) {Система пуста};
    \node[below, font=\footnotesize] at (8.5, -0.4) {Система пуста};

\end{tikzpicture}
\end{center}

\textbf{2. Интегральное тождество.}
Рассмотрим площадь фигуры, заключенной между графиками $A(t)$ и $D(t)$. С одной стороны, интегрируя по вертикали, эта площадь равна $\int_{0}^{T} (A(t) - D(t)) dt$. С другой стороны, разбивая фигуру на горизонтальные полосы, длина каждой $k$-й полосы в точности равна времени пребывания $k$-й заявки в системе $W^{(k)}$. Следовательно, справедливо тождество суммарного времени:
$$ \int_{0}^{T} N(t) dt = \sum_{k=1}^{N} W^{(k)}. $$

\textbf{3. Асимптотический переход.}
Поделим обе части полученного тождества на $T$ и умножим правую часть на $\frac{N}{N}$:
$$ \underbrace{\frac{1}{T} \int_{0}^{T} N(t) dt}_{L(T)} = \underbrace{\frac{N}{T}}_{\lambda(T)} \cdot \underbrace{\left( \frac{1}{N} \sum_{k=1}^{N} W^{(k)} \right)}_{W(N)}. $$
Устремляя $T \to \infty$ (а следовательно, и $N \to \infty$), мы немедленно получаем искомое соотношение:
$$ \lim_{T \to \infty} L(T) = \lim_{T \to \infty} \lambda(T) \cdot \lim_{N \to \infty} W(N) \implies L = \lambda \cdot W. $$

\textbf{4. Обобщение на произвольную дисциплину очереди.}
Покажем, что предположение о порядке выхода заявок (FCFS) не является существенным. Пусть $A^{(k)}$ — момент поступления, а $D^{(k)}$ — момент выхода $k$-й заявки (при этом индексы выхода могут не быть монотонными). Факт нахождения $k$-й заявки в системе в момент времени $t$ описывается индикаторной функцией $\mathbb{I}(A^{(k)} \le t < D^{(k)})$. Тогда общее число заявок:
$$ N(t) = \sum_{k=1}^{N} \mathbb{I}(A^{(k)} \le t < D^{(k)}). $$
Интегрируя $N(t)$ по времени на отрезке $[0, T]$ и меняя местами сумму и интеграл, получаем:
$$ \int_{0}^{T} N(t) dt = \int_{0}^{T} \left( \sum_{k=1}^{N} \mathbb{I}(A^{(k)} \le t < D^{(k)}) \right) dt = \sum_{k=1}^{N} \int_{0}^{T} \mathbb{I}(A^{(k)} \le t < D^{(k)}) dt. $$
Поскольку интеграл от индикатора на отрезке $[0, T]$ равен длине интервала пребывания заявки в системе (с учетом того, что система пуста в момент $T$), он в точности равен $W^{(k)}$. Таким образом, интегральное тождество $\int_0^T N(t) dt = \sum_{k=1}^N W^{(k)}$ инвариантно относительно дисциплины обслуживания.

\textbf{5. Краевые эффекты.}
Откажемся от требования пустоты системы в момент $T$. Если в момент $T$ в системе остаются необработанные заявки, то площадь фигуры даст погрешность $\varepsilon(T)$, равную суммарному времени, которое эти «зависшие» заявки провели в системе до момента $T$. Однако в стационарном (эргодическом) режиме очередь не растет бесконечно. Следовательно, остаточная площадь $\varepsilon(T)$ ограничена величиной порядка $O(1)$, и при передельном переходе вклад краевых эффектов исчезает:
$$ \lim_{T \to \infty} \frac{\varepsilon(T)}{T} = 0. $$
Теорема полностью доказана.
\end{tproof}